\chapter{Results and Discussion}\label{ch:results}
\section{Results and Analysis}
\subsection{Run Integrity and Coverage}
The analyzed run is \code{three-topics-dev-v1-fixed-01}. It contains all 225 expected samples and 225 score artifacts. The recorded report audit checks sample--score fingerprints, question and qrels alignment, recomputed retrieval metrics, aggregate breakdowns, paired statistics and Holm corrections, finite values, and quality/latency CSV exports. The numeric inputs in this report are exported from those frozen artifacts; they do not come from the earlier smoke test.

Every configuration completes 43 of 45 questions and returns insufficient evidence for GD014 and KHCN013. No inference execution is recorded as failed or timed out. Across configurations, this is 215 completed responses and ten insufficient-evidence outcomes. The 95.56\% completion rate describes workflow coverage, not answer accuracy.

\begin{figure}[H]\centering
\begin{tikzpicture}[node distance=6mm]
\node[box,text width=11cm] (a) {45 questions per method: retrieval and context scores available};
\node[box,text width=11cm,below=of a] (b) {43 completed answers + 2 insufficient-evidence outcomes};
\node[box,text width=11cm,below=of b] (c) {43 answer-relevancy scores; 2 completed answers have claim-extraction errors};
\node[box,text width=11cm,below=of c] (d) {41 valid faithfulness / factual-score pairs per method};
\draw[arrow] (a)--(b);\draw[arrow] (b)--(c);\draw[arrow] (c)--(d);
\end{tikzpicture}
\caption{Metric availability per configuration. Different denominators prevent a single uniform ``accuracy'' interpretation.}\label{fig:availability}
\end{figure}

The two extraction-error questions are GT009 and GD015 in every configuration, giving ten judge-error samples overall. The judge output is truncated at its configured 2,048-token limit during answer-claim extraction. These are scoring errors on completed inference artifacts. They are not reclassified as inference failures, and missing claim-based scores are excluded rather than replaced with zero. Context metrics remain available for all 45 questions and answer relevancy for 43.

\subsection{Retriever Performance}
Table~\ref{tab:retrieval} reports macro means over all 45 questions for each method. Recall@5 is already high for the hybrid baseline, while B1, B3 and B4 obtain slightly higher observed recall. The small nDCG differences indicate changes near the top of an already strong ranking on this development set.

\begin{table}[H]\centering\small
\caption{Retrieval quality on 45 development questions. All values use binary leaf qrels.}\label{tab:retrieval}
\begin{tabular}{lrrrrrr}\toprule
\textbf{Method} & \textbf{R@5} & \textbf{P@5} & \textbf{MRR@5} & \textbf{MAP@5} & \textbf{nDCG@5} & \textbf{R@20}\\\midrule
\input data/retrieval-table.tex
\bottomrule\end{tabular}
\note{Precision@5 has a dataset-label ceiling of 0.2267; the reference labels are not exhaustive judgments of every useful passage.}
\end{table}

\begin{figure}[H]\centering
\begin{tikzpicture}
\begin{axis}[width=.94\textwidth,height=5.8cm,ymin=0,ymax=1.08,symbolic x coords={B0,B1,B2,B3,B4},xtick=data,ylabel={Macro mean},legend style={at={(.5,-.20)},anchor=north,legend columns=3}]
\addplot+[mark=*,color=navy,thick] table[x=method,y=recall5,col sep=comma]{data/metrics.csv};\addlegendentry{Recall@5}
\addplot+[mark=square*,color=teal,thick] table[x=method,y=ndcg5,col sep=comma]{data/metrics.csv};\addlegendentry{nDCG@5}
\addplot+[mark=triangle*,color=orange,thick] table[x=method,y=precision5,col sep=comma]{data/metrics.csv};\addlegendentry{Precision@5}
\addplot[black,dashed] coordinates {(B0,0.22666667) (B4,0.22666667)};\addlegendentry{P@5 label ceiling}
\end{axis}
\end{tikzpicture}
\caption{High labeled-evidence recall coexists with a low numeric precision ceiling. Lines connect configurations for readability, not a continuous parameter.}\label{fig:retrieval}
\end{figure}

B0 obtains Recall@5 of 0.9667 and P@5 of 0.2133. The reranked variants reach 0.9778 and 0.2178. Against the qrels ceiling, those precision values represent most of the possible top-five relevant-label count. It would be incorrect to diagnose a 78\% answering error rate from $1-P@5$, because the denominator is five retrieved ranks and most questions label only one leaf.

B4 reaches Recall@20 of 1.0000, yet still does not complete the two rejected questions. Retrieval recall and answer sufficiency therefore diverge in an observable way: finding the labeled leaf IDs somewhere in a ranked list is not equivalent to retaining all required context or obtaining a correct gate decision.

\subsection{Packed Context and Answer Quality}
\begin{table}[H]\centering\small
\caption{Context and answer metrics from the current offline judge. These are provisional measurements.}\label{tab:quality}
\begin{tabular}{lrrrrr}\toprule
\textbf{Method} & \textbf{Context P} & \textbf{Context R} & \textbf{Faithfulness} & \textbf{Factual F1} & \textbf{Answer rel.}\\\midrule
\input data/quality-table.tex
\bottomrule\end{tabular}
\note{Valid questions per method: Context P/R $n=45$; faithfulness and factual F1 $n=41$; answer relevancy $n=43$. Higher is better within each metric, but their units and denominators differ.}
\end{table}

Context Precision ranges from 0.9854 to 0.9963, while Context Recall ranges from 0.7697 to 0.8305. This is consistent with a context collection that the judge finds useful early in its order but not necessarily complete against all reference claims. High rank-aware context precision should not be interpreted as evidence that every clause needed by the question is present.

B2 has the largest observed Context Recall and Factual F1 means, whereas B3 and B4 have the largest faithfulness and answer-relevancy means. B1 increases retrieval Recall@5 relative to B0 but has a lower judged factual F1. These descriptive differences show why ranking and answer scores must both be inspected. They do not prove which component caused a particular judgment, especially given the judge errors discussed below.

\begin{figure}[H]\centering
\begin{tikzpicture}
\begin{axis}[width=.94\textwidth,height=5.8cm,ymin=0,ymax=1.05,symbolic x coords={B0,B1,B2,B3,B4},xtick=data,ylabel={Offline judged macro mean},legend style={at={(.5,-.20)},anchor=north,legend columns=2}]
\addplot+[color=navy,mark=*,thick] table[x=method,y=context_precision,col sep=comma]{data/metrics.csv};\addlegendentry{Context Precision ($n=45$)}
\addplot+[color=teal,mark=square*,thick] table[x=method,y=context_recall,col sep=comma]{data/metrics.csv};\addlegendentry{Context Recall ($n=45$)}
\addplot+[color=orange,mark=triangle*,thick] table[x=method,y=faithfulness,col sep=comma]{data/metrics.csv};\addlegendentry{Faithfulness ($n=41$)}
\addplot+[color=purple,mark=diamond*,thick] table[x=method,y=factual_f1,col sep=comma]{data/metrics.csv};\addlegendentry{Factual F1 ($n=41$)}
\end{axis}
\end{tikzpicture}
\caption{Stage-specific quality profiles. Judge calibration is required before treating these answer scores as reliable legal-correctness estimates.}\label{fig:quality}
\end{figure}

The additional claim-precision, claim-recall and strict-yield results appear in Table~\ref{tab:factual}. Near-zero strict yield follows from requiring two imperfect judge measures to equal one simultaneously. It should not be described as evidence that nearly all answers are wrong.

\begin{table}[H]\centering\small
\caption{Reference-based claim scores and strict correct-answer yield.}\label{tab:factual}
\begin{tabular}{lrrrr}\toprule
\textbf{Method} & \textbf{Factual P} & \textbf{Factual R} & \textbf{Factual F1} & \textbf{Strict yield}\\\midrule
\input data/factual-table.tex
\bottomrule\end{tabular}
\note{Claim metrics use 41 valid completed answers. Yield uses 43 valid outcomes, including the two insufficient-evidence zeros and excluding two judge-error answers.}
\end{table}

\subsection{Paired Uncertainty}
No reported test reaches a Holm-adjusted $p<0.05$. Figure~\ref{fig:ci} displays the paired nDCG@5 deltas relative to B0. For B3--B0, the mean difference is 0.0088 with a pointwise 95\% interval of $[-0.0124,0.0347]$ and adjusted $p=1.000$. This does not substantiate a retrieval advantage over the baseline.

\begin{figure}[H]\centering
\begin{tikzpicture}
\begin{axis}[width=.90\textwidth,height=5.6cm,xmin=-.025,xmax=.045,ymin=.5,ymax=4.5,ytick={1,2,3,4},yticklabels={B1--B0,B2--B0,B3--B0,B4--B0},xlabel={Paired change in nDCG@5},ylabel={Comparison},xtick={-.02,0,.02,.04}]
\draw[dashed,gray] (axis cs:0,.5)--(axis cs:0,4.5);
\addplot[navy,thick] coordinates {(-0.013576412829684055,1) (0.03379717417095429,1)};\addplot[only marks,mark=*,navy] coordinates {(0.007854115164926642,1)};
\addplot[navy,thick] coordinates {(0.0,2) (0.015086292368727782,2)};\addplot[only marks,mark=*,navy] coordinates {(0.005028764122909261,2)};
\addplot[navy,thick] coordinates {(-0.012433143731305425,3) (0.03466085314926534,3)};\addplot[only marks,mark=*,navy] coordinates {(0.008798681249596381,3)};
\addplot[navy,thick] coordinates {(-0.010250744063789714,4) (0.03637465048070691,4)};\addplot[only marks,mark=*,navy] coordinates {(0.010582663712969991,4)};
\end{axis}
\end{tikzpicture}
\caption{Point estimates and pointwise 95\% paired bootstrap intervals ($n=45$, 2,000 resamples). All four Holm-adjusted tests have $p=1.000$; intervals are not simultaneous family-adjusted intervals.}\label{fig:ci}
\end{figure}

B2's factual F1 delta relative to B0 is 0.0157 on 41 paired valid questions, with pointwise interval $[0.0000,0.0423]$ and adjusted $p=1.000$. A non-negative pointwise interval is not sufficient to declare a winner when the reported multiple-comparison procedure does not reject the null. Conversely, failing to find significance is not proof of equivalence. The experiment is small, development-tuned and judged by a model that still needs calibration.

\subsection{Inference Latency}
\begin{table}[H]\centering\small
\caption{Full-pipeline wall-clock latency in seconds. Preparation and offline scoring are excluded.}\label{tab:latency}
\begin{tabular}{lrrrr}\toprule
\textbf{Method} & \textbf{All mean} & \textbf{Completed mean} & \textbf{Completed p50} & \textbf{Completed p95}\\\midrule
\input data/latency-table.tex
\bottomrule\end{tabular}
\note{All mean uses $n=45$; completed-only statistics use $n=43$ per method. Percentiles use the benchmark's linear interpolation convention. These are measurements on one incompletely described local machine.}
\end{table}

B0 has the lowest observed all-outcome mean at 21.76 seconds. B1, B2 and B3 average 23.04, 22.13 and 23.27 seconds. B4 averages 24.28 seconds, approximately 11.6\% above B0 in this run. These values describe the recorded local execution; one repeat and incomplete hardware notes do not support a portable latency guarantee.

\begin{figure}[H]\centering
\begin{tikzpicture}
\begin{axis}[width=.94\textwidth,height=6cm,xlabel={Full-pipeline seconds per question},ylabel={Cumulative fraction},xmin=8,xmax=50,ymin=0,ymax=1.03,legend style={at={(.98,.02)},anchor=south east,legend columns=1}]
\addplot+[const plot,mark=none,thick,color=navy] coordinates {(14.62896,0.022222) (17.69378,0.044444) (18.30951,0.066667) (18.42239,0.088889) (18.42276,0.111111) (18.91455,0.133333) (18.97225,0.155556) (19.03853,0.177778) (19.32617,0.200000) (19.41245,0.222222) (19.44464,0.244444) (19.70799,0.266667) (19.96633,0.288889) (20.39853,0.311111) (20.59578,0.333333) (20.69512,0.355556) (20.90056,0.377778) (20.96515,0.400000) (21.24384,0.422222) (21.31699,0.444444) (21.32714,0.466667) (21.35033,0.488889) (21.38807,0.511111) (21.39630,0.533333) (21.51391,0.555556) (21.54695,0.577778) (21.73645,0.600000) (21.81106,0.622222) (21.81498,0.644444) (22.01691,0.666667) (22.08185,0.688889) (22.16797,0.711111) (22.39391,0.733333) (22.48783,0.755556) (22.64062,0.777778) (22.71109,0.800000) (23.09403,0.822222) (23.95613,0.844444) (24.73112,0.866667) (25.24563,0.888889) (25.59825,0.911111) (25.66639,0.933333) (27.86674,0.955556) (28.91788,0.977778) (35.21316,1.000000)};\addlegendentry{B0}
\addplot+[const plot,mark=none,thick,color=orange] coordinates {(18.00896,0.022222) (18.02045,0.044444) (18.79274,0.066667) (19.01234,0.088889) (19.10552,0.111111) (19.20042,0.133333) (19.47592,0.155556) (19.61007,0.177778) (19.87920,0.200000) (19.90587,0.222222) (20.46993,0.244444) (20.56661,0.266667) (20.63214,0.288889) (20.97408,0.311111) (21.00165,0.333333) (21.26225,0.355556) (21.37513,0.377778) (21.40349,0.400000) (21.47482,0.422222) (21.70685,0.444444) (21.71491,0.466667) (21.84037,0.488889) (21.85571,0.511111) (21.86332,0.533333) (21.93131,0.555556) (21.99525,0.577778) (22.18571,0.600000) (22.24726,0.622222) (22.34969,0.644444) (22.47496,0.666667) (22.73062,0.688889) (22.75921,0.711111) (23.28784,0.733333) (23.44697,0.755556) (23.93002,0.777778) (24.27742,0.800000) (24.59594,0.822222) (27.60061,0.844444) (28.01006,0.866667) (28.20952,0.888889) (28.95568,0.911111) (31.28641,0.933333) (32.54856,0.955556) (36.09953,0.977778) (36.76732,1.000000)};\addlegendentry{B1}
\addplot+[const plot,mark=none,thick,color=teal] coordinates {(14.79799,0.022222) (18.06504,0.044444) (18.35934,0.066667) (18.89028,0.088889) (19.09015,0.111111) (19.31079,0.133333) (19.36718,0.155556) (19.46785,0.177778) (19.56523,0.200000) (19.59844,0.222222) (20.05396,0.244444) (20.17323,0.266667) (20.74340,0.288889) (20.84071,0.311111) (21.10539,0.333333) (21.15332,0.355556) (21.20587,0.377778) (21.31273,0.400000) (21.36305,0.422222) (21.47835,0.444444) (21.50089,0.466667) (21.70639,0.488889) (21.72075,0.511111) (21.76018,0.533333) (21.92734,0.555556) (22.00708,0.577778) (22.02686,0.600000) (22.03213,0.622222) (22.10391,0.644444) (22.18094,0.666667) (22.33578,0.688889) (22.55730,0.711111) (22.65885,0.733333) (22.75466,0.755556) (22.84601,0.777778) (23.26947,0.800000) (23.88551,0.822222) (24.56177,0.844444) (25.22399,0.866667) (25.41468,0.888889) (25.61806,0.911111) (27.40345,0.933333) (28.39433,0.955556) (29.31457,0.977778) (34.90097,1.000000)};\addlegendentry{B2}
\addplot+[const plot,mark=none,thick,color=purple] coordinates {(17.51746,0.022222) (18.05961,0.044444) (18.75466,0.066667) (18.91916,0.088889) (19.10492,0.111111) (19.49815,0.133333) (19.78149,0.155556) (19.91681,0.177778) (20.07825,0.200000) (20.46571,0.222222) (20.60035,0.244444) (20.72331,0.266667) (20.98306,0.288889) (20.99993,0.311111) (21.12316,0.333333) (21.36529,0.355556) (21.42580,0.377778) (21.83107,0.400000) (21.87237,0.422222) (21.92079,0.444444) (22.11154,0.466667) (22.13868,0.488889) (22.18139,0.511111) (22.22881,0.533333) (22.41022,0.555556) (22.48717,0.577778) (22.52742,0.600000) (22.74425,0.622222) (23.04078,0.644444) (23.54213,0.666667) (23.59150,0.688889) (23.90539,0.711111) (24.24572,0.733333) (24.27050,0.755556) (24.38636,0.777778) (24.42240,0.800000) (24.78766,0.822222) (26.73255,0.844444) (28.89036,0.866667) (28.94163,0.888889) (29.20195,0.911111) (30.30871,0.933333) (31.46514,0.955556) (33.90570,0.977778) (37.83801,1.000000)};\addlegendentry{B3}
\addplot+[const plot,mark=none,thick,color=black] coordinates {(18.27036,0.022222) (19.08115,0.044444) (19.17524,0.066667) (19.47434,0.088889) (19.49497,0.111111) (19.65679,0.133333) (19.89945,0.155556) (20.20322,0.177778) (20.42927,0.200000) (20.52799,0.222222) (20.95022,0.244444) (20.97689,0.266667) (21.21580,0.288889) (21.27582,0.311111) (21.31589,0.333333) (21.44229,0.355556) (21.58998,0.377778) (21.77219,0.400000) (21.81050,0.422222) (22.02597,0.444444) (22.04416,0.466667) (22.07525,0.488889) (22.28953,0.511111) (22.37997,0.533333) (22.44849,0.555556) (22.54052,0.577778) (22.56475,0.600000) (22.81097,0.622222) (23.43452,0.644444) (23.56372,0.666667) (24.01909,0.688889) (24.24887,0.711111) (24.51711,0.733333) (24.91296,0.755556) (25.20645,0.777778) (28.05776,0.800000) (28.79881,0.822222) (29.16195,0.844444) (30.03539,0.866667) (30.41481,0.888889) (32.04573,0.911111) (32.85949,0.933333) (37.09557,0.955556) (39.60965,0.977778) (44.89602,1.000000)};\addlegendentry{B4}
\end{axis}
\end{tikzpicture}
\caption{Empirical latency distributions over all 45 outcomes per method, including insufficient-evidence cases. Curves show observations rather than fitted distributions.}\label{fig:latency}
\end{figure}

The completed-only means for B3 and B4 are very similar. Their difference in all-outcome mean is driven by extra work on the rejected questions: B4's two insufficient-evidence outcomes average 42.25 seconds, compared with 18.78 seconds for B3. It is therefore important to include those outcomes when describing the operational cost of feedback. Reporting completed-only speed would hide the additional unsuccessful retrieval rounds.

\subsection{Feedback and Graph Observations}
In B4, 43 questions stop after one retrieval round and two run all three rounds. The total is 49 rounds, compared with 45 for each single-round method. All 45 final answer objects are identical between B3 and B4. Feedback changes some deeper ranking results, including Recall@20, but recovers no additional completed answer in this pilot.

The graph branch is enabled in B2, B3 and B4. Inspection of stored retrieval rounds finds zero retained relation lists in the final ranked retrieval results for all methods, and zero explicit document-relation blocks in final context text. This does not mean the prepared graph has no relationships or that the enabled branch could not affect candidate ranks. It does mean that the pilot provides no evidence of answering with an explicit retained graph-relation block. Graph-assisted candidate retrieval and demonstrated graph-grounded multi-hop reasoning are different claims.

Mean packed-context sizes are approximately 1,489 tokens for B0, 1,498 for B1, 1,476 for B2, 1,498 for B3 and 1,504 for B4. These averages are below the configured context ceiling, but do not establish that every needed provision was available or selected. A larger budget alone cannot fix a wrong candidate or an overly demanding gate.

\section{Discussion}
\subsection{Why Completion Is Not Accuracy}
The absence of inference failures is a useful engineering result. The revised contracts can finish the complete B0--B4 workflow and preserve inspectable artifacts. Nevertheless, the remaining insufficient-evidence cases and quality-judge issues prevent a claim of 95.56\% legal accuracy. Completed responses still require evidence-level review, and a fallback can be appropriate when the requested facts are not retained.

The two rejected questions illustrate different bottlenecks. For GD014, the packed material misses the general scope provision of the Education Law and emphasizes narrower adult-education and policy provisions. Rejecting that context can be justified by the missing scope. For KHCN013, the retained material includes the three appendices and the clause concerning proposals to update the list, but the online evaluator demands additional procedural deadlines and documents beyond the requested detail. This suggests a possible false-negative sufficiency judgment. It is not evidence that the generator should always be allowed to answer after a failed gate.

\subsection{Why the Offline Scores Need Calibration}
The stored judgments permit direct counterexamples to a literal reading of faithfulness. In B0/GD001, the answer and the retained scope provision both enumerate the education system, educational institutions, teachers, learners, state management, and associated rights and responsibilities. The judge nevertheless reports faithfulness of 0.3333 while factual F1 is 1.0000. Its reasons demand detailed provisions or confirmation of current validity that the scope statement does not assert.

In B0/GT002, the answer states the applicable organizations and individuals for inland-waterway activity, with a treaty condition. The factual judge rejects part of the answer using a reason about other transport modes that the response does not claim. These cases demonstrate semantic judging errors. They do not show that every low score is erroneous, nor establish an adjusted true score.

Claim extraction is also a source of error: removing the relation ``the law regulates'' can turn a supported statement about scope into a fragmented statement that the entailment judge interprets differently. GT009's extraction truncation occurs even though its final response is only about 318 characters. Increasing the token limit might reduce a particular failure, but would not by itself demonstrate stable or valid claim decomposition. The required next step is calibrated extraction and entailment against human-reviewed examples, followed by rescoring the same frozen answers.

\begin{table}[H]\centering\small
\caption{Observed issues separated by pipeline or measurement stage.}\label{tab:issues}
\begin{tabularx}{\textwidth}{P{2.3cm}P{3cm}Y}\toprule
\textbf{Case} & \textbf{Stage} & \textbf{Evidence and interpretation}\\\midrule
GD014 & Retrieval / context & General scope not retained; missing evidence can explain fallback.\\
KHCN013 & Online sufficiency & Evaluator requests details beyond the asked scope; possible false negative.\\
GT009, GD015 & Offline extraction & Truncation in all methods; claim-based scores unavailable.\\
B0/GD001 & Offline entailment & Same scope information receives low faithfulness but perfect factual F1.\\
B0/GT002 & Offline factual judgment & Rejection reason introduces transport modes absent from the response.\\\bottomrule
\end{tabularx}
\end{table}

\subsection{Threats to Validity}
\textbf{Construct validity.} Binary leaf labels, rank-aware context precision and claim-based quality scores represent different constructs. None is a complete measure of legal correctness. Incomplete alternative-evidence labeling constrains the interpretation of precision, and model judgments introduce prompt- and model-dependent error.

\textbf{Internal validity.} The ablation shares corpus, models, filters and budgets, which supports comparison. However, a component changes the downstream evidence seen by common stages, so aggregate changes cannot be attributed to one isolated mechanism without trace analysis. Near-ceiling retrieval and a single repeat reduce sensitivity to small gains. Development fixes and question clarifications also prevent an unbiased final-test interpretation.

\textbf{External validity.} Only three topics and 45 answerable questions are included. Explicit document and title cues make some questions easier to constrain than broad real-user requests. Few questions are comparisons or multi-document tasks, and the run includes no held-out unanswerable or adversarial temporal cases. Generalization to open-ended legal consultation or graph-dependent reasoning remains untested.

\textbf{Measurement reliability.} Some judge outputs are invalid and others semantically inconsistent. Local hardware and LLM quantization records are incomplete. A technically correct aggregation can faithfully summarize unreliable judgments; arithmetic consistency and semantic validity must therefore be audited separately.

\clearpage
\section{Recommendations}
The first priority is to stabilize claim extraction and calibrate the offline judge on a stratified, independently reviewed subset. Reviewers should preserve conditions, exceptions and the original question's scope. Agreement, error categories, and examples should be documented before introducing a new judge identity and rescoring the existing inference artifacts.

The second priority is a reviewed dataset with a development/test boundary established before further tuning. It should include alternative supporting provisions where relevant, more multi-document questions requiring real graph links, and unanswerable cases. Topic balance alone is insufficient; task difficulty, evidence multiplicity, document age and ambiguity should also be characterized. Negative controls should test whether related but insufficient provisions are correctly rejected.

The third priority is a controlled repeated experiment with complete CPU/GPU/RAM, model digest, quantization, database and load records. Paired comparisons should be planned in advance with a small set of primary metrics. B0 remains the simplest reference configuration for continued development because it has the lowest observed mean latency and strong retrieval here. This is an engineering choice for the next experiment, not a statistically established quality winner.

For graph research, instrument the number of graph-only relevant hits, retained relation evidence, and genuinely graph-dependent questions. For feedback, distinguish recovery of missing evidence from repeated rejection of already adequate context. These measurements would explain \emph{when} an addition helps, which is more informative than enabling it by default because it increases architectural complexity.
